\begin{tabularx}{\textwidth}{L{34mm}X X}
\toprule
\textbf{Vezérlési stílus} & \textbf{Alapgondolat} & \textbf{Jellemző változatok / alkalmazás} \\
\midrule
Központosított vezérlés & Egy kitüntetett vezérlő alrendszer felel a többi alrendszer indításáért, leállításáért és koordinálásáért. & Hívás-visszatérés modell szekvenciális rendszerekben; kezelőmodell szekvenciális vagy konkurens folyamatok koordinálására. \\
Hívás-visszatérés modell & A vezérlés a felsőbb szintű rutintól alprogramhívásokkal jut az alsóbb szintekre, majd visszatér a hívás helyére. & Könnyen elemezhető, de merev; normálistól eltérő működésnél és erősen eseményfüggő rendszereknél kényelmetlen lehet. \\
Kezelőmodell & Egy rendszerkezelő komponens állapotváltozók és események alapján ciklikusan indítja, leállítja és koordinálja a folyamatokat. & Konkurens rendszereknél is használható; gyakran eseményciklus jellegű vezérlőként jelenik meg. \\
Eseményalapú vezérlés & A vezérlés nincs egyetlen komponensbe zárva; az alrendszerek eseményekre reagálnak. & Eseményszóró modell, grafikus felületek, szerkesztők, lazábban csatolt rendszerek. \\
Megszakítás-vezérelt modell & Külső megszakításokat egy megszakításkezelő észlel, majd a megfelelő komponenshez irányít. & Valós idejű és beágyazott rendszerekben jellemző, ahol gyors külső eseményekre kell reagálni. \\
\bottomrule
\end{tabularx}
